\documentclass[11pt,a4paper]{article}
\usepackage[T1]{fontenc}
\usepackage{lmodern}
\pdfmapfile{+lm.map}
\pdfmapfile{+symbols.map}
\pdfmapfile{+cm.map}
\pdfmapfile{+cmextra.map}
\usepackage[margin=25mm]{geometry}
\usepackage{amsmath,amssymb,amsthm,booktabs,longtable,array,enumitem}
\usepackage[dvipsnames]{xcolor}
\usepackage{microtype}
\usepackage{fancyhdr}
\usepackage[breaklinks,colorlinks,linkcolor=MidnightBlue,urlcolor=MidnightBlue,citecolor=MidnightBlue]{hyperref}
\usepackage{listings}
\newtheorem{theorem}{Theorem}[section]
\newtheorem{lemma}[theorem]{Lemma}
\newtheorem{proposition}[theorem]{Proposition}
\newtheorem{corollary}[theorem]{Corollary}
\newcommand{\PA}{\mathrm{PA}}
\newcommand{\Con}{\operatorname{Con}}
\newcommand{\Halt}{\operatorname{Halt}}
\newcommand{\Nonhalt}{\operatorname{Nonhalt}}
\newcommand{\code}[1]{\texttt{#1}}
\setlength{\headheight}{14pt}
\setlength{\emergencystretch}{2em}
\pagestyle{fancy}
\fancyhf{}
\fancyhead[L]{\small Con(PA) machine comparison}
\fancyhead[R]{\small 13 November 2026}
\fancyfoot[C]{\thepage}
\setlist{itemsep=3pt,topsep=5pt}
\lstset{basicstyle=\ttfamily\small,columns=fullflexible,breaklines=true,frame=single,rulecolor=\color{black!20}}
\title{\textbf{Con(PA) Turing machines:\\comparison, verified intermediate constructions,\\and the remaining state-reduction obstruction}}
\author{Research and reproducibility report}
\date{13 November 2026}
\begin{document}
\maketitle
\begin{abstract}
The requested machine with fewer than 196 working states was not obtained.
This report does not certify a 195-state machine. It compares the supplied
208-, 209-, and 196-state designs at pinned revisions, accounts for the
different 208-state presentations, and records the new work in a form that
can be checked and continued. The principal complete intermediate combines
a cached-head-only Shift test with an MP result stored in that cache. It
has 196 binary working states, 3,788 source instruction occurrences, an exact
57,922-step blank startup, and 4,928 uniform source-instruction contracts.
A written proof supplies its PA-provable consistency calibration; it is
distinct from a new Lean kernel-checked endpoint. A repaired 157-node
source layout preserves the old 3,812-instruction control graph, but the
verified physical allocation has 207 states. A direct-negation specialization
reduces the source to 3,556 instructions and has a verified 197-state
realization. An even-tag formula recoding
justifies a simpler parity acceptance condition, but no sub-196-state
implementation or complete parity startup is certified here.
\end{abstract}
\tableofcontents
\newpage

\section{What was requested, and what is established}
The goal is a deterministic binary Turing machine $M$ on a single two-way
infinite tape, started on the all-zero tape, with fewer than 196 working
states and
\[
          \PA\vdash\Nonhalt(M)\longleftrightarrow\Con(\PA).
\]
The distinguished halting label is not a working state. Every working state
has exactly two transition rows, one for each scanned symbol; a transition
writes a bit and moves exactly one cell left or right. Changing the alphabet,
using a preloaded tape, or counting source DAG nodes as machine states
would not meet this goal.

The requested strict inequality remains open in this investigation. The
accepted complete intermediate in this packet is denoted
$M_{196}^{\mathrm{combined}}$. Its transition table is
\path{machines/h_only_mp_h.tm}. The file contains every transition, not a
patch requiring an unspecified ancestor. Its control labels are
$0,\ldots,195$, its start label is 0, and its terminal label is 196.

\begin{center}\small
\begin{tabular}{@{}lrrl@{}}
\toprule
Object & Working states & Source sites & Status in this packet\\
\midrule
Pinned PA196 &196&3812&Reference table and uniform checks\\
Combined Shift/MP &196&3788&Complete written proof and checks\\
Direct negation variant &197&3556&Equivalent arithmetic macro\\
Repaired 157-node source &207&3812&Same source CFG; verified realization\\
Parity control, reset length 8 &198&3801&Runtime component only\\
Requested construction &$\le195$&---&Not obtained\\
\bottomrule
\end{tabular}
\end{center}

There are three different kinds of evidence throughout this report. Exact
finite checks establish table counts, hashes, local rewrites, finite source
bindings, and the literal startup configuration. Uniform source-site checks
use symbolic tape runs and the unchanged register-engine lemmas to cover
arbitrary natural register values. Written arithmetic arguments establish
the changed commands and the proof compiler. None of these is described
as a completed new Lean formalization of the modified machine.

\section{Pinned inputs and comparison method}
\subsection{Revisions and archive coverage}
All comparisons use the following revisions, rather than a moving main
branch:
\begin{center}\small
\begin{tabular}{@{}ll@{}}
\toprule
Repository/version & Commit\\
\midrule
BusyBeaver, PA196 &\texttt{407cdc356f6b586635e71ac92be394255f4bfea4}\\
BusyBeaver, PA209 &\texttt{69d98241a721ead20b4f0b17df9062950b40fbfe}\\
ProveIt, incoming and ConciseIPRA &\texttt{dc523e7b06ffb414093c90aa035500ca283e5f75}\\
\bottomrule
\end{tabular}
\end{center}
The current incoming directory contains 18 top-level ZIP archives. The
inventory in \path{provenance/current_archive_inventory.json} records their
exact hashes, byte sizes, entry counts, and report paths. The GitHub blob
identities were checked when materializing these archives. For every
selected 208-state comparison table, the packet additionally records every
matching current ZIP member in
\path{provenance/comparison_input_verification.json}. This prevents an old
scratch copy from silently substituting for the current supplied input.

The seven newer reports in the incoming directory contain prior reductions,
negative experiments, and repeated copies of earlier tables. They are part
of the literature being compared. Their already-recorded unsuccessful
195-state candidates, HA-strength modification of ConciseIPRA, and written
internal-converse arguments are not claimed as discoveries of this report.
The full source archives total well over 100 MiB; the deliverable retains
selected exact tables, relevant source, hashes, provenance, and notices
instead of recursively duplicating those archives.

\subsection{Canonicalizing machines without confusing syntax and behavior}
The canonical comparison traverses the rooted transition graph in breadth
first order, taking the read-0 edge before the read-1 edge. It records, for
each edge, the written bit, move direction, and canonical successor label,
with a distinguished terminal marker. Equality of these encodings is rooted,
action-preserving graph isomorphism. It is stronger than equal state count
and weaker than deciding whether two different machines have the same
halting behavior. It is a finite structural test, not a solution to machine
equivalence.

\path{comparison/compare_designs.py} checks input hashes and complete tables,
performs this canonicalization, and checks the extra local helper blocks.
It runs both normally and under \texttt{python3 -O}; its essential comparisons
use explicit exceptions rather than removable Python assertions.

\section{How the supplied designs differ}
\subsection{The 208-state family has more than two classes}
Nine designated presentations in the five original U208 packets collapse
to two rooted graph classes, called A and B in the comparison receipt.
However, \path{pa208_research_2026-10-29.zip} supplies 18 additional helper
layouts: nine based on A and nine based on B. Each has a distinct rooted
transition graph, and none is isomorphic to A or B. The correct designated
comparison therefore has
\[
       9+18=27\text{ presentations},\qquad 2+18=20\text{ graph classes}.
\]
Saying that all supplied 208-state machines form only two classes would
omit those 18 layouts. Copies embedded in later research packets do not
create further designated designs merely by being copied.

\begin{center}\small
\begin{tabular}{@{}lcc@{}}
\toprule
Original packet (all dated 2026-10-29)&Primary class&Alternative class\\
\midrule
\texttt{U208\_ConPA}&A&---\\
\texttt{U208\_ConPA\_Conditional\_Context}&A&B\\
\texttt{U208\_ConPA\_Entry\_Context}&A&B\\
\texttt{U208\_ConPA\_History}&A&B\\
\texttt{u208\_research}&B&A (plain variant)\\
\bottomrule
\end{tabular}
\end{center}

The common original 208-state reduction identifies the historical native
controls 172 and 211. Their entry histories expose different predecessor
bits, so the merged control can recover the required choice from the tape.
One presentation keeps 172 and removes 211; another naming keeps 211.
These are presentations of the same saving, not independent savings that
can be stacked to obtain 207 states. The four whole-tape local blocks have
lengths 7, 1, 3, and 1 in the representative audit. With the historical
state labels, the stuttering-time relation is
\[
 T_{208}=T_{209}+6N(172,0)+2N(211,0),\qquad
 T_{209}\le T_{208}\le7T_{209}.
\]
Here $N(q,b)$ counts the old trace's visits to the indicated read case.
The relation concerns the designated U209 parent and its U208 descendant.
It must not be applied by label name to James's different PA209 table.

The extra 18 helper layouts identify another pair and add a helper control,
so their net state count remains 208. Their finite helper blocks are replayed
in the comparison script. They give useful different sharing arrangements,
but no additional state saving by themselves.
\input{chapters/helper_layouts.tex}

\subsection{The two 209-state ancestors must be distinguished}
The supplied U208 packets descend from their own U209 parents. Those parents
are not literally James's PA209 machine at the historical commit. They share
an earlier architecture, but source bracketings, selected sharing, labels,
and startup traces differ. A direct subtraction of numerical labels across
the two branches is consequently invalid.

\begin{center}
\begin{tabular}{@{}lrrrl@{}}
\toprule
Design & States & Rows & Source sites & Reported/checked startup\\
\midrule
Original U208 A/B &208&416&3909&48140\\
James PA209 &209&418&3948&58210\\
James PA196 &196&392&3812&58154\\
Combined196 &196&392&3788&57922\\
\bottomrule
\end{tabular}
\end{center}
The U208 source has 5,083 resolved graph edges and 5,084 outcome contracts
in the inherited reports. PA196 has 4,959 outcome contracts, counting both
zero and positive cases of each decrement. The startup counts are exact
reaches of specific clean frames, not useful-run benchmarks or evidence
of nonhalting. The PA196 and Combined196 startups are replayed in this
packet; the older counts are attributed to their pinned construction
certificates and comparative sources.

\subsection{PA196 is a logical and compilation change}
PA196 is not obtained by simply applying all the old U208 identifications
to PA209. In native numbering, its comparison with James's PA209 removes
75 labels, introduces 62 new 1300-series labels, and keeps 134 labels.
Among those retained labels, 265 transition rows are identical. The only
three retained-row changes are
\[
\begin{array}{c|c|c}
\text{native row}&\text{PA209}&\text{PA196}\\\hline
(101,1)&1R115&1R1367\\
(164,0)&0L311&0R77\\
(190,1)&1L308&1R99.
\end{array}
\]
Both old controls 172 and 211 are already absent. Reapplying their old
identification cannot reduce the 196-state table.

The physical allocation consists of 36 engine states, 13 leaf-routing
states, and 147 reader states representing 159 distinct source nodes.
The AX body is a 201-item macro word. Its 68 distinct new reader nodes are
represented by 62 physical controls through six contextual folds. The
overall source contains further sharing outside AX.

Its constructors use Cantor pairing
\[
 p(a,b)=\frac{(a+b)(a+b+1)}2+b,
\quad I(a,b)=4p(a,b),\quad M(a,b)=4p(a,b)+1,
\]
\[
 U(a,B)=4p(B,a)+2,\qquad E(a,b)=4p(a,b)+3.
\]
The transposition in $U$ is intentional. Code 0 acts as truth and code 1
is the chosen contradiction. The logical changes include moving equality
vacuity into the guarded separation bundle, removing the old separate
$\rho_7$ role, and changing some head orientations and their ordering.
These changes shorten the AX word and improve physical sharing. The exact
definitions, rather than a verbal summary, are the specification; see
\texttt{PA196/Commands/Code.lean} in the pinned source \cite{bb196}.

\section{The logical strength and the meaning of ``provable in PA''}
Let $\mathcal C$ be the numerical calculus searched by PA196. Four separate
claims are needed for the requested calibration:
\begin{enumerate}
\item PA refutations translate to $\mathcal C$ refutations.
\item $\mathcal C$ refutations translate syntactically to PA refutations.
\item A $\mathcal C$ refutation yields a halting command word that is usable
by the exhaustive scheduler, whatever store previous trials have left.
\item The literal blank-tape machine and the source program simulate one
another for halting, with an exact startup and positive finite steps.
\end{enumerate}
For an internal biconditional, these transformations and their correctness
must be available within PA, not merely as true external assertions.

The upstream file \path{PA196/Independence.lean} includes an internal
theorem named \path{pa_proves_nonhalting_consistency}. Its statement is
\[
          \PA\vdash\Nonhalt(M_{196})\longrightarrow\Con(\PA).
\]
It also exports nonprovability and two-sided independence, using the
separate external soundness/model development. This should not be described
as an exported Lean proof of the displayed biconditional. In particular,
external soundness of a calculus does not by itself supply the reverse
implication inside PA.

The reverse internal direction has a written syntactic argument in the
earlier incoming research, reproduced with exact provenance in
Appendix~\ref{sec:reverse-consistency}. Decode formulas by a primitive recursive
arithmetization, interpret membership as binary-bit membership of natural
numbers, and compile each calculus axiom and inference into a PA proof of
its arithmetic translation. Finite-comprehension instances are justified
by arithmetic induction on the finite bound. This constructs proofs of
individual translated instances; it does not install a truth predicate
for arbitrary arithmetic formulas. A numeric-calculus proof of code 1
then yields a PA proof of a fixed contradiction. Correctness is established
by induction on the coded derivation.

The full calibration is consequently a written mathematical theorem with
explicit finite and uniform components, while the absent endpoint remains
a formalization task. The new combined construction below retains the same
calculus and adapts its command compiler. It does not claim that a new
Lean endpoint has been built or that the upstream project was rebuilt in
this run. Selected upstream definitions and proof statements are included
for review; complete dependency reproduction follows the pinned repository's
own instructions.

\section{The new complete intermediate: Combined196}
\subsection{The six physical changes}
Starting with the dense PA196 table, change precisely these six rows:
\begin{center}
\begin{tabular}{@{}rrccc@{}}
\toprule
State&Read&Write&Move&Next\\
\midrule
115&1&1&R&108\\
37&0&0&R&11\\
37&1&1&R&101\\
127&0&0&R&133\\
127&1&1&R&20\\
133&0&0&R&120\\
\bottomrule
\end{tabular}
\end{center}
The first change moves the MP success result into the cached head. The
remaining changes implement the Shift variant that tests the cached head
alone. The existing single halting row, $(63,1)\mapsto(1,R,196)$, remains.
The full table's SHA-256 is
\begin{quote}\small
\nolinkurl{bd74754d75e7b8ba97fa7bfc104374de33ae603fc183ae7f08f053ab8c9ea592}.
\end{quote}

The two component changes were already present separately in prior research;
the new result here is their simultaneous realization and the completed
command compiler for the resulting behavior. Combining the patches without
updating that compiler would leave a logical gap. In particular, Pop becomes
partial at a saved top of 1, and the MP compiler needs a dummy zero beneath
its premises to preserve the underlying stack.

The combined source has 159 distinct internal DAG nodes and 3,788 leaf
occurrences, compared with 159 and 3,812 in the original. Thus 24 fewer
source occurrences do not imply any fewer physical states. The longest
source address is 39 bits. All 196 working controls are graph reachable.

\subsection{Exact startup and uniform runtime certificate}
The table reaches root control 75 with head at absolute coordinate $-2$
after exactly 57,922 literal transitions. The whole tape is checked against
the clean frame at source site 25. Relative to the head, that frame consists
of a left marker at $-1$, the path \texttt{1000000}, enough zeros to put the
file marker at 48, and the encoded unary register file. The physical values
are
\[
 [0,0,0,0,0,0,2,0,0]\mathbin{+\!+}[1]^{26}.
\]
Under the physical register permutation
$[6,2,0,4,3,7,8,5,1]$, the logical candidate register is 2 and all other live
logical registers are zero. The test checks the complete finite support
of the tape, not just the first several cells or the fact that control 75
was visited.

The symbolic verifier checks 4,928 instruction outcomes over every natural
register file and every additional clean gap $n\ge0$, with total gap at
least 48. A symbolic zero run has length $n+c$; the remaining file is an
opaque arbitrary register-file object. Scan lemmas and four register-effect
summaries stand for unbounded loops. Their hypotheses are checked against
the unchanged engine rows before use. The maximum symbolic trace length
is 524 verifier steps. This number is a proof-checking statistic, not a
bound on literal execution for large register values.

The union of the literal startup rows and the rows visited in symbolic
source execution covers every one of the 392 transition rows. No state or
row can be dropped by the simple unused-row argument under this certificate.
The arithmetic source-macro proof and command-level calibration are given
in Sections~\ref{sec:combined-source-macros} and
\ref{sec:combined-command-compiler}.

\input{chapters/combined_source_macros.tex}
\input{chapters/combined_commands.tex}

\section{Direct negation: a shorter source with a 197-state realization}
\label{sec:direct-negation}
A further arithmetic specialization replaces every occurrence of the AX
macro sequence $\mathsf P\,\mathsf i\,\mathsf C$. In the inherited notation,
these macros push the current top, increment the new top, then combine the
two values into an implication code. Their composite constructs
$N(F)=I(F,1)=4p(F,1)$ while preserving the underlying stack. There are eight
such occurrences in the 201-item AX word.

For arbitrary natural $t,x,y$, the existing \texttt{node46} pairing core
has the uniform arithmetic contract
\[
          (T,X,Y)=(t,x,y)\longmapsto(0,0,y+p(x,t)),
\]
with $W$ and the other framed registers unchanged. To verify it, first
transfer $T$ into $X$ and $Y$. This adds $t$ to $Y$ and leaves $X=x+t$.
The subsequent passes add the successive values $x+t,x+t-1,\ldots,1$
to $Y$, decreasing that loop rank after each pass. Thus the total addition
is $t+(x+t)(x+t+1)/2=p(x,t)$. This is an induction on natural loop ranks,
not a finite sample test. The existing \texttt{node63} then transfers $Y$
into the initially zero $T$ four units at a time, giving $(T,Y)=(4Y,0)$.
Their composition is \texttt{node64}.

Consequently the closed macro
\[
        N_{\mathrm{direct}}=
          (\operatorname{inc}Y,(\operatorname{inc}Y,\texttt{node64}))
\]
sends $(T,X,Y)=(F,0,0)$ to $(4p(F,1),0,0)$, since
\[
        p(0,F)+2=p(F,1).
\]
It preserves arbitrary $W$. The usual zero-scratch invariant holds at AX
macro boundaries. Replacing all eight occurrences therefore leaves every
AX output code and its stack order unchanged. All of Combined196's logical
command proofs transfer directly to this variant.

The exact word/source binding is checked by
\path{verification/check_negation.py}. The AX word shrinks from 201 to
185 macro items and the full source from 3,788 to 3,556 instructions.
Compiling the selected grammar gives 202 states before five simultaneous
compatible folds, and 197 afterwards. The table
\path{machines/direct_negation197.tm} passes all 4,616 uniform source
outcomes and reaches the exact clean seed after 57,922 literal steps.
Its labels are sparse; its terminal label remains 196 and is excluded from
the count of 197 working labels. Its SHA-256 is
\begin{quote}\small
\nolinkurl{832e067d50599a0ed9ead8e19d621e6491d04faebd821e7018d3aa6708786fa6}.
\end{quote}
Two other equal-value macro arrangements were examined: transfer $F$ into
$X$, set $T=1$, and use \texttt{node64}; or use \texttt{node64} first and
then add eight to $T$ with existing increment blocks. The latter has worse
address geometry in the tested layouts. The accepted 197-state table is
a new complete arithmetic specialization, but it costs one more physical
state than PA196 and two more than the requested maximum.

\section{Why the repaired 157-node source still costs 207 states}
Earlier incoming research supplies a 157-node source tree with the same
3,812-instruction ordinal control graph. Its original physical attempt
failed at a transient zero-run case. This investigation reassociated the
source to shorten that problematic increment suffix, then repaired a
continue-to-depth-3 leaf whose path had changed from \texttt{101011} to
\texttt{1010011}. A subsequent exact-CFG rotation restored a workable jump
context without losing the 157-node count.

The selected source is
\path{research/source157/source157_tailrepair_5.json}. Its resolved
instruction graph is exactly equal to the original PA196 graph, including
the register named by every primitive instruction and both outcomes of
every decrement. Its maximum address depth is 36 and its maximum unmodified
zero run is 6. Its emitted physical table passes all 4,959 symbolic
instruction contracts and reaches the exact clean seed after 58,140 steps.

The physical table has 207 working states. A trace-based census identifies
153 reader controls representing the 157 source nodes, plus 54 nonreader
controls. The original table has 147 readers for 159 source nodes plus 49
nonreaders. Existing helper paths keep several old source fragments alive;
the generic compiler also loses contextual identifications. Exact action
bisimulation leaves all 207 controls distinct. The tested finite-context
direct detours and compatible two-cell folds provide no further saving for
this selected allocation. These are scoped search results, not a lower
bound on every possible allocation of this source.

This example is a useful failure diagnosis: a correct, smaller DAG and a
correct, larger TM can coexist. A claim of a 195-state machine based only
on $157<159$ would leave out the physical decoder, route, and reset costs.

\input{chapters/parity_recoding.tex}
\input{chapters/membership_acceptance.tex}

\section{Physical experiments for the parity architecture}
The mathematical recoding in Section~\ref{sec:parity-recoding} is a proposed
architecture, not the encoding currently present in the parity-control
table. The control experiment changes the Shift prelude to leave
$T=h$, $H=0$, and $X=t+h$; removes the final increment of $X$; and uses
the source guard
\[
       (\operatorname{dec}X,
          ((\operatorname{dec}X,\operatorname{continueTo}2),\operatorname{halt})).
\]
It consumes $X$ in pairs and halts on odd inputs. The source contains
3,801 sites and 4,944 outcome contracts.

\begin{center}
\begin{tabular}{@{}lrrl@{}}
\toprule
Return detector & States & Failing outcomes & Meaning\\
\midrule
Original seven-zero detector&197&38&19 source sites fail\\
Shortened six-zero detector&196&105&Invalid runtime certificate\\
Extended eight-zero detector&198&0&All 4944 local contracts pass\\
\bottomrule
\end{tabular}
\end{center}
The original runtime has intermediate returns through the root decoder on
19 source sites. They are harmless with its original root but not with the
new parity root. The eight-zero detector repairs these local contracts at
the cost of one state. It does not discharge startup: the old exact seed
was not reached in the explicit two-million-step test. Earlier clean frames
were observed with different register contents, so this is an unresolved
bootstrap obligation, not a proof that the table can never work.

Moreover, these parity tables still use the old four-residue AX/MP
constructors. They have not integrated the even-tag encoding needed by the
logical argument. The table is deliberately placed under
\path{research/parity_and_engine}, not under the complete machines.

The base-6 and base-8 grammar searches charged for the new multiplier
macro and equality tag. Their best tested layout estimates did not improve
on the original physical count; some were around 200 even before a fully
validated parity guard and bootstrap. These estimates are neither certified
machine counts nor optimality results. The parity theorem remains useful
because it identifies a valid simpler acceptance condition and precisely
states the encoding change needed to support it.

\section{ConciseIPRA: useful ideas and the strength barrier}
The current ConciseIPRA README \cite{ipra} describes substantially more than
a small checker file. It includes raw one-term transport, direct progress
rules, a seven-rule single-introduction presentation, clause-based five-
and four-family forms, and three-family forms such as Use/Intro/Name and
Use/Intro/Open on full normal forms. It distinguishes Lean proofs of
calculus equivalence from Python measurements of checker byte length.

These results provide useful methods: derive redundant rules, use a
normal form that exposes common code, and separate logical equivalence from
the cost of its implementation. They do not automatically produce a
binary Con(PA) machine with the same number of rules, bytes, or source
nodes. A normalizer, substitution machinery, proof enumerator, and input
encoding all have a physical state cost.

There is also a logical distinction. A machine searching for an iPRA
refutation naturally calibrates to iPRA consistency. Its equivalence to
PA consistency requires an additional proof-theoretic argument; one cannot
replace the target theory by name. The prior incoming packet
\path{pa196_research_2026-11-13.zip} already records the HA-strength
generalization obtained by allowing compound induction progress formulas,
with the corresponding PA consistency calibration. That is relevant prior
work. No smaller binary implementation of that strengthened checker was
completed here.

\section{Rejected shortcuts and exact search limits}
\begin{enumerate}
\item \textbf{Combining the 208-state merges.} The original packets remove
the same native 172/211 distinction. Renaming the surviving state does not
make a second independent reduction.
\item \textbf{Deleting a reset state.} The nominal 195-state six-zero
version fails the first source instruction's uniform runtime contract.
The clean source address has six zeros, and a transient runtime marker
creates the misleading interior boundary. This refutes the proposed
simulation certificate; it does not decide the blank-run behavior of
every 195-state machine.
\item \textbf{Accepting all odd codes in the original encoding.}
Equality has residue 3. The attempted clone interpretation fails because
equality can distinguish objects that the proposed membership relation
cannot separate. Even-tag recoding is needed for the parity argument.
\item \textbf{Keeping an old compiler after changing commands.} The new
Pop can halt on a saved top of 1, and the MP suffix consumes an extra stack
entry unless it is supplied a dummy zero. Both facts are explicitly handled
by the compiler in this report.
\item \textbf{Removing native state 198.} All 780 tested one-row redirects
of its sole incoming edge failed the 72 selected register-file cases, even
when the allowed spare tail was relaxed. This rules out that finite family,
not every possible redesign of the carry engine.
\item \textbf{Five-control synthesis of a six-control carry block.}
Unrestricted bounded synthesis timed out. A restricted version preserving
the original nonzero scan rows was unsatisfiable on the selected sample
contracts and bounds. No engine minimality conclusion follows.
\item \textbf{Fresh-label and weak-head proposals.} Enlarged fresh labels
have an unproved low-label vacuity/refutation-translation obligation. They
have not been disproved merely by that gap. Countermodels to restricted
implicational fragments do not establish incompleteness of the entire
quantifier-and-gate calculus.
\end{enumerate}

New finite-window helper searches were also run on PA196, the combined
table, an earlier structurally different 196-state bracketing, and its
combined variant. The helper may perform arbitrary allowed bit writes and
left/right moves within the stated finite window before rejoining the
old transition endpoint. No compatible net-saving shared helper or direct
detour merge was found in those searches. This is a restricted absence of
candidates; it is not evidence that 196 is minimal.

A separate exact-CFG physical-rotation search, keeping the AX region fixed,
exhausted its connected component of 1,000 layouts at depth 16. All had
196 states after the tested duplicate-row and predecessor-bit folds.
This finite component is defined by the allowed rotation moves, not by
all possible source bracketings or all equivalent Turing machines.

For Combined196, an independent unary/list symbolic checker also rederived
its engine summaries and passed all 4,928 cases, recording 1,204,510
symbolic events. Intersecting its source contexts with an independent
blank-window closure left 2,214 control/window pairs. An engine-inclusive
screen then considered one deleted state and up to two arbitrary changed
rows, requiring the original transition endpoint within 16 child steps
and within the five cells at offsets $-2,\ldots,2$. It tested 38,220 ordered
state deletions, 118,490,624 first-row assignments, and 576,694,720
second-row assignments. All 635 representative passes failed on the full
admitted context set. The two narrower deletion templates also had no
survivors. The exact context hash, programs, scope, and counts are under
\path{research/engine_screen}. This is a negative result for the stated
local simulation templates, not for every possible 195-state machine.

\section{Reproduction and certificate boundaries}
From the packet directory run:
\begin{lstlisting}
python3 verify.py
python3 verify.py --extended
pdflatex -interaction=nonstopmode -halt-on-error report.tex
pdflatex -interaction=nonstopmode -halt-on-error report.tex
\end{lstlisting}
The optional \texttt{python3 build\_report.py} wrapper performs the two
LaTeX passes and can build a local format when a TeX installation supplies
the packages but has no prebuilt format.
The first command checks the pinned PA196 and Combined196 tables, exact
startups, all their symbolic source outcomes, the canonical comparison,
the 64 symbolic command-algebra checks, and the finite source-macro binding.
The second also checks the source-CFG identity and 207-state realization
of the repaired 157-node tree, together with the 197-state direct-negation
variant. Python needs only its standard library.
The command and binding supplements are also executed under optimization
to ensure their essential conditions are not lost with Python assertions.

The scripts output machine-readable receipts under \path{receipts}.
\path{SHA256SUMS.json} records the packaged file hashes. The upstream Lean
source is included selectively as a reference and proof dependency; this
packet's Python verification is not a substitute for rebuilding the
upstream Lean project. The written changed-command proof is new work,
while its AX, scheduler, pairing, register-engine, and proof-translation
dependencies are explicitly inherited from the pinned sources and prior
written arguments.

The deliverable records a failed strict state-reduction goal and useful
verified intermediate results. It supplies no certified transition table
with fewer than 196 working states. A future success must provide that
table, its actual physical census, an exact blank startup, a uniform
runtime proof, and the PA-internal proof transformations together.

\appendix
\clearpage
\section*{Attribution for the reproduced proof appendix}
The following section is reproduced unchanged from
\path{pa_tm_research_2026-11-13.zip}, member
\path{pa_tm_research_2026-11-13/report/pa_reverse_translation.tex}.
Its exact archive and member hashes are in
\path{reference/prior_proofs/provenance.json}. It is prior work supplied
among the requested inputs. Its implementation obligations and conditional
transfer statement are retained as written. Sections 6 and 7 of the present
report supply the changed source and command arguments for Combined196;
the new machine still has no newly exported Lean endpoint. The appendix's
encoded negations use the fixed PA theorem $\neg B(x_0,x_0)$, so their
translation has the intended falsity behavior rather than an arbitrary
propositional constant.
\input{reference/prior_proofs/pa_reverse_translation.tex}
\begin{thebibliography}{9}
\raggedright
\bibitem{bb196}
James Chang, \emph{BusyBeaver: PA196}, pinned source and paper,
commit \texttt{407cdc356f6b586635e71ac92be394255f4bfea4}.
\url{https://github.com/jamesyc/BusyBeaver/tree/407cdc356f6b586635e71ac92be394255f4bfea4}.
\bibitem{bb209}
James Chang, \emph{BusyBeaver: PA209}, historical source and paper,
commit \texttt{69d98241a721ead20b4f0b17df9062950b40fbfe}.
\url{https://github.com/jamesyc/BusyBeaver/tree/69d98241a721ead20b4f0b17df9062950b40fbfe}.
\bibitem{incoming}
Vladimir Reshetnikov, \emph{ProveIt, docs/incoming}, 18 pinned research
archives, commit \texttt{dc523e7b06ffb414093c90aa035500ca283e5f75}.
\url{https://github.com/VladimirReshetnikov/ProveIt/tree/dc523e7b06ffb414093c90aa035500ca283e5f75/docs/incoming}.
\bibitem{ipra}
Vladimir Reshetnikov, \emph{ConciseIPRA}, current README and linked
calculus developments at the same ProveIt commit.
\url{https://github.com/VladimirReshetnikov/ProveIt/tree/dc523e7b06ffb414093c90aa035500ca283e5f75/Logic/PrimitiveRecursiveArithmetic/ConciseIPRA}.
\bibitem{prior}
Earlier research packets included in \cite{incoming}, especially
\path{pa196_research_2026-11-13.zip},
\path{pa208_research_2026-10-29.zip}, and the prior comparison/reduction
reports. Exact current archive and selected-member hashes are listed in
this packet's provenance directory. Prior proof arguments and code are
reused with their original notices preserved.
\end{thebibliography}
\end{document}
